固有多項式$|A- \lambda E|= 0$を計算する。
1)より行列式が0であるため、$\lambda= 0$は固有値の1つ。また行列$A$のトレースと主小行列式の和を用いて計算することもできる。
\begin{itemize}
    \item $\tr{A}= 4+ 3- 4= 3$
    \item $\beign{vmatrix}3& -6\\ 2& -4\end{vmatrix}+ \begin{vmatrix}4& 2\\ -6& -4\end{vmatrix}+ \begin{vmatrix}4& -1\\ -6& 3\end{vmatrix}= 2$
    \item |A|= 0
\end{itemize}
    したがって固有方程式は
\begin{equation*}
    \lambda^3- 3\lambda\2+ 2\lambda= \lambda(\lambda- 1)(\lambda- 2)= 0
\end{equation*}
    固有値は$\lambda= 0, 1, 2$

    \item 各固有値に対応する固有ベクトルを求める。
\begin{enumerate_alpha}
    \item $\lambda= 0$：(2)より固有ベクトル$\bm{v}_{\lambda= 0}= \begin{pmatrix}0& 2\\ 1\end{pmatrix}$
    \item $\lambda_ 2$：$(A- E)\bm{x}= 0$を解けばよい。
\begin{equation*}
    \begin{pmatrix}
        3& -1& 2\\
        -6& 2& -6\\
        -6& 2& -5
    \end{pmatrix}
    \begin{pmatrix}x_1\\ x_2\\ x_3\end{pmatrix}
    = \begin{pmatrix}0\\ 0\\ 0\end{pmatrix}
\end{equation*}
    これを解くと$x_3= 0, $



    が成り立つ。相似な行列はトレースが等しいため、$\tr{A}= \tr{B}$である。ここで元の行列$A$のトレースは$\tr{A}= 3$であるため、$\tr{B}= 3$とならなければならない。
    しかし過程より$B$のすべての成分は偶数であるため、対角成分もすべて偶数。したがって$\tr{B}= b_{11}+ b_{22}+ \b_{33}$も偶数とならなければならない。これは$\tr{B}= 3$に矛盾する。したがって成分がすべて偶数となるような表現行列をもつ基底$F$は存在しない。


    \item 半角の公式$\sin^2{\theta}= \frac{1- \cos{2\theta}}{2}$を用いる。
    \begin{align*}
        \int_0^{2\pi} \sin^2{x}dx
        &= \int_0^{2\pi} \frac{1- \cos{2x}}{2}dx\\
        &= \left[\frac{x}{2}- \frac{\sin{2x}{4}}\right]_0^{2\pi}\\
        &= \left(\frac{2\pi}{2}- 0\right)- (0- 0)= \pi
    \end{align*}
    \item $r= \sin{u}$とおく。
\begin{equation}
    dr= \cos{u}du,\ \     
    \begin{tablar}{l|l}
        r& 0\to 2\pi\\
        u& 0\to \frac{\pi}{2}
    \end{tablar}
\end{equation}

\begin{align*}
    \int_0^1\r^2\sqrt{1- r^2}dr
    &= \int_0^{\frac{\pi}{2}}\sin^2{u}\sqrt{1- \sin^2{u}}(\cos{u})du\\
    &= \int_0^{\frac{\pi}{2}}\sin^2{u}\cos^2{u} du\\
    &= \int\frac{1}{4}\sin^2{(2u)}du\\
    &= \frac{1}{4}\int_0^{\frac{\pi}{2}}\frac{1- \cos{(4u)}}{2}du\\
    &= \frac{1}{8}\left[u- \frac{\sin{(4u)}}{4}\right]_0^{\fra{\pi}{2}}\\
    &= \frac{\pi}{16}
\end{align*}

    \item 与えられた変数変換$\begin{cases}x= r\cos{\theta},\\ y= 2r\sin{\theta}\end{cases}$について、ヤコビアンを求める。
\begin{align*}
    \frac{\partial(x, y)}{\partial (r, \theta)}
    = \begin{vmatrix}\frac{\partial x}{\partial r}& \frac{\partial x}{\partial \theta}\\ \frac{\partial y}{\partial r}& \frac{\partial y}{\partial \theta}\\
    &= \begin{vmatrix}\cos{\theta}& -r\sin{\theta}\\ 2\sin{\theta}& 2r\sin{\theta}\end{vmatrix}\\
    &= 2r\cos^2{\theta}+ 2r\sin^2{\theta}= 2r
\end{align*}
